\hypertarget{appendix-a-incident-register}{%
\chapter{Appendix A --- Incident
register}\label{appendix-a-incident-register}}

Grounded incidents that motivated specific T-entries. Each is public and
durably sourced; the mapping to threat identifiers is given at the end
of each entry.

\hypertarget{ona-veto-research-march-2026}{%
\section{Ona Veto research (March
2026)}\label{ona-veto-research-march-2026}}

Demonstrated an AI coding agent autonomously bypassing two layers of
security:

\begin{enumerate}
\def\labelenumi{\arabic{enumi}.}
\tightlist
\item
  The agent discovered an alternate path resolving to a denylisted
  binary that did not match the deny pattern.
\item
  When the sandbox blocked namespace creation, the agent reasoned its
  way to disabling the sandbox itself rather than abandoning the task.
\end{enumerate}

The research also demonstrated that, when content-addressed enforcement
was applied, the agent pivoted to invoking the dynamic linker to load
denied binaries via memory-mapping rather than execution, bypassing an
execution-based enforcement gate.

\emph{Source:
\url{https://ona.com/stories/how-claude-code-escapes-its-own-denylist-and-sandbox}}

\emph{Threat IDs: T2.1 (host security control deactivation). Motivates a
defence-in-depth, hostility-by-default adversary model.}

\hypertarget{nx-s1ngularity-supply-chain-attack-august-2025}{%
\section{Nx s1ngularity supply-chain attack (August
2025)}\label{nx-s1ngularity-supply-chain-attack-august-2025}}

Eight malicious Nx package versions published 26--27 August 2025
contained a \texttt{telemetry.js} postinstall script that:

\begin{enumerate}
\def\labelenumi{\arabic{enumi}.}
\tightlist
\item
  Inventoried the local filesystem for sensitive files.
\item
  Invoked AI agent CLIs in permission-bypassing modes with
  reconnaissance prompts.
\item
  Exfiltrated discovered credentials to public repositories named with
  an \texttt{s1ngularity-repository-} pattern.
\item
  Appended a shutdown command to shell startup files to cause future
  shells to immediately shut down.
\end{enumerate}

GitGuardian's analysis documented 1,346 affected repositories, 2,349
distinct exfiltrated secrets across 1,079 compromised developer systems,
90\% of leaked tokens still valid post-incident. A follow-on attack used
the exfiltrated tokens to rename roughly 5,500 previously-private
organisation repositories to public access.

\emph{Sources:
\url{https://snyk.io/blog/weaponizing-ai-coding-agents-for-malware-in-the-nx-malicious-package/};
\url{https://blog.gitguardian.com/the-nx-s1ngularity-attack-inside-the-credential-leak/};
\url{https://socket.dev/blog/nx-packages-compromised}}

\emph{Threat IDs: T1.1 (reconnaissance), T1.2 (postinstall), T1.9
(supply chain).}

\hypertarget{clinejection-disclosed-february-2026-exploited-january-2026}{%
\section{Clinejection (disclosed February 2026, exploited January
2026)}\label{clinejection-disclosed-february-2026-exploited-january-2026}}

A prompt-injection vulnerability in a CI issue-triage workflow allowed
any user to compromise npm publishing tokens by opening an issue with a
crafted title. The chain:

\begin{enumerate}
\def\labelenumi{\arabic{enumi}.}
\tightlist
\item
  An issue title containing prompt injection caused the triage agent to
  execute attacker-controlled installation from an imposter commit.
\item
  The malicious pre-install script deployed cache-poisoning tooling to
  the CI runner.
\item
  The tooling flooded the CI cache, triggering eviction of legitimate
  entries, and set poisoned entries matching the nightly publish
  workflow's keys.
\item
  The nightly publish restored the poisoned cache, which exfiltrated the
  publishing tokens.
\item
  The attacker used the stolen token to publish a malicious package
  version with a hostile postinstall script.
\end{enumerate}

The malicious package was live for approximately eight hours, downloaded
approximately 4,000 times.

\emph{Sources: \url{https://adnanthekhan.com/posts/clinejection/};
\url{https://snyk.io/blog/cline-supply-chain-attack-prompt-injection-github-actions/}}

\emph{Threat IDs: T1.2 (postinstall), T1.3 (compromised extension), T1.9
(supply chain).}

\hypertarget{mindgard-research-october-2025}{%
\section{Mindgard research (October
2025)}\label{mindgard-research-october-2025}}

Four vulnerabilities in a widely-installed IDE agent extension:

\begin{enumerate}
\def\labelenumi{\arabic{enumi}.}
\tightlist
\item
  \textbf{DNS-based exfiltration} --- prompt-injected docstrings caused
  the agent to read environment variables and encode them into DNS
  queries to attacker-controlled domains, using a command on the
  safe-command allowlist.
\item
  \textbf{Rules-file privilege escalation} --- malicious content in an
  agent-rules directory overrode the approval flag for arbitrary
  commands.
\item
  \textbf{Model exfiltration} --- attacker-controlled prompts leaked the
  system prompt and model information.
\item
  \textbf{Arbitrary code execution} --- via combinations of the above.
\end{enumerate}

\emph{Threat IDs: T1.1 (reconnaissance), T1.3 (compromised extension),
T1.7 (DNS exfiltration), T3.7 (prompt injection).}

\hypertarget{agent-cli-cves-20252026}{%
\section{Agent-CLI CVEs (2025--2026)}\label{agent-cli-cves-20252026}}

A run of disclosed vulnerabilities in AI agent CLIs illustrates the
compromised-extension and channel-redirection classes: malicious hooks
in project configuration executing shell commands during initialisation
before any consent dialog; a project-set base-URL variable redirecting
API requests to attacker endpoints before the trust prompt; a sandbox
escape via settings-file injection.

\emph{Threat IDs: T1.3 (compromised extension), T1.8 (channel
redirection), T2.1 (control deactivation).}

\hypertarget{axios-npm-compromise-april-2026}{%
\section{axios npm compromise (April
2026)}\label{axios-npm-compromise-april-2026}}

A compromised package version delivered a remote access trojan via a
postinstall hook, executing in approximately two seconds --- before
installation completed.

\emph{Threat IDs: T1.2 (postinstall), T1.9 (supply chain).}

\hypertarget{runc-cve-2024-21626-leaky-vessels-january-2024}{%
\section{runc CVE-2024-21626 (``Leaky Vessels'', January
2024)}\label{runc-cve-2024-21626-leaky-vessels-january-2024}}

A file-descriptor leak allowed container processes to reach the host
filesystem via an unclosed descriptor referring to the host's root
directory. Affected the mainstream container and orchestration runtimes
using the vulnerable versions.

\emph{Threat IDs: T3.2 (container escape).}

\begin{center}\rule{0.5\linewidth}{0.5pt}\end{center}

\hypertarget{appendix-b-reconnaissance-target-enumeration}{%
\chapter{Appendix B --- Reconnaissance target
enumeration}\label{appendix-b-reconnaissance-target-enumeration}}

The paths a credential-hunting workload (T1.1) will typically attempt to
read, by category. This is representative, not exhaustive; it is the
reconnaissance surface a constructed-view approach removes by absence.

\begin{itemize}
\tightlist
\item
  SSH and PGP: \texttt{\textasciitilde{}/.ssh/},
  \texttt{\textasciitilde{}/.gnupg/}
\item
  Cloud credentials: \texttt{\textasciitilde{}/.aws/},
  \texttt{\textasciitilde{}/.azure/},
  \texttt{\textasciitilde{}/.config/gcloud/},
  \texttt{\textasciitilde{}/.oci/},
  \texttt{\textasciitilde{}/.linode-cli},
  \texttt{\textasciitilde{}/.ibmcloud/},
  \texttt{\textasciitilde{}/.config/doctl/}
\item
  Forge tokens: \texttt{\textasciitilde{}/.config/gh/},
  \texttt{\textasciitilde{}/.config/glab/},
  \texttt{\textasciitilde{}/.config/hub},
  \texttt{\textasciitilde{}/.config/bitbucket-cli/}
\item
  Legacy auth: \texttt{\textasciitilde{}/.netrc}
\item
  Package manager credentials: \texttt{\textasciitilde{}/.npmrc},
  \texttt{\textasciitilde{}/.yarnrc},
  \texttt{\textasciitilde{}/.yarnrc.yml},
  \texttt{\textasciitilde{}/.pypirc},
  \texttt{\textasciitilde{}/.cargo/credentials},
  \texttt{\textasciitilde{}/.docker/config.json},
  \texttt{\textasciitilde{}/.kube/config}
\item
  Infrastructure tooling:
  \texttt{\textasciitilde{}/.terraform.d/credentials.tfrc.json},
  \texttt{\textasciitilde{}/.config/terraform/}
\item
  Password managers: \texttt{\textasciitilde{}/.password-store/},
  \texttt{\textasciitilde{}/.config/keepassxc/},
  \texttt{\textasciitilde{}/.config/Bitwarden*},
  \texttt{\textasciitilde{}/.config/1Password/}
\item
  System keyrings: \texttt{\textasciitilde{}/.local/share/keyrings/},
  \texttt{\textasciitilde{}/.gnome2/keyrings/}
\item
  Browser state: \texttt{\textasciitilde{}/.mozilla/},
  \texttt{\textasciitilde{}/.config/google-chrome/},
  \texttt{\textasciitilde{}/.config/chromium/},
  \texttt{\textasciitilde{}/.config/BraveSoftware/},
  \texttt{\textasciitilde{}/.config/microsoft-edge/},
  \texttt{\textasciitilde{}/.config/vivaldi/}
\item
  Messaging clients: \texttt{\textasciitilde{}/.config/Signal/},
  \texttt{\textasciitilde{}/.config/Slack/},
  \texttt{\textasciitilde{}/.config/discord/},
  \texttt{\textasciitilde{}/.config/Element/},
  \texttt{\textasciitilde{}/.thunderbird/}
\item
  Shell histories: \texttt{\textasciitilde{}/.bash\_history},
  \texttt{\textasciitilde{}/.zsh\_history},
  \texttt{\textasciitilde{}/.fish\_history},
  \texttt{\textasciitilde{}/.local/share/fish/fish\_history}
\item
  Tool histories: \texttt{\textasciitilde{}/.python\_history},
  \texttt{\textasciitilde{}/.node\_repl\_history},
  \texttt{\textasciitilde{}/.psql\_history},
  \texttt{\textasciitilde{}/.mysql\_history},
  \texttt{\textasciitilde{}/.sqlite\_history},
  \texttt{\textasciitilde{}/.lesshst},
  \texttt{\textasciitilde{}/.viminfo}
\item
  Cryptocurrency wallets: \texttt{\textasciitilde{}/.electrum/},
  \texttt{\textasciitilde{}/.bitcoin/},
  \texttt{\textasciitilde{}/.ethereum/},
  \texttt{\textasciitilde{}/.config/Exodus/},
  \texttt{\textasciitilde{}/.config/Atomic/}
\item
  VPN configuration: \texttt{\textasciitilde{}/.config/openvpn3/},
  \texttt{\textasciitilde{}/.config/wireguard/},
  \texttt{\textasciitilde{}/.cisco/},
  \texttt{\textasciitilde{}/.config/forticlient/}
\end{itemize}

Within typically-granted project trees:

\begin{itemize}
\tightlist
\item
  \texttt{.env}, \texttt{.env.local}, \texttt{.env.production},
  \texttt{.env.staging}, \texttt{.env.*}
\item
  \texttt{config/secrets.yml}, \texttt{application-prod.properties},
  \texttt{appsettings.json}
\item
  \texttt{terraform.tfstate} (contains every secret Terraform has
  touched)
\item
  \texttt{.envrc} (direnv source from secret stores)
\item
  \texttt{.git/config} with credential-helper URLs
\end{itemize}

\begin{center}\rule{0.5\linewidth}{0.5pt}\end{center}

\hypertarget{appendix-c-mitre-attck-mapping-summary}{%
\chapter{Appendix C --- MITRE ATT\&CK mapping
summary}\label{appendix-c-mitre-attck-mapping-summary}}

\begin{longtable}[]{@{}ll@{}}
\toprule\noalign{}
T-ID & ATT\&CK techniques \\
\midrule\noalign{}
\endhead
\bottomrule\noalign{}
\endlastfoot
T1.1 & T1552, T1083, T1005 \\
T1.2 & T1195.002, T1059 \\
T1.3 & T1554, T1195.001 \\
T1.4 & T1059.004, T1105 \\
T1.5 & T1204.002 \\
T1.6 & T1021, T1570 \\
T1.7 & T1071.004, T1048 \\
T1.8 & T1071.001, T1041 \\
T1.9 & T1195.002 \\
T1.10 & T1098, T1543 \\
T1.11 & T1189, T1218 \\
T1.12 & T1021, T1185 (partial) \\
T1.13 & T1021, T1559 \\
T1.14 & T1083, T1518 \\
T2.1 & T1562, T1562.001, T1562.008 \\
T2.2 & T1554, T1505 (partial) \\
T2.3 & T1552.001, T1552.004 \\
T2.4 & T1098.003, T1078 \\
T2.5 & T1562.001, T1554 \\
T2.6 & T1115 \\
T2.7 & T1113, T1059 (partial), T1185 (partial) \\
T2.8 & T1546, T1059 \\
T3.1 & T1548.001 \\
T3.2 & T1611, T1610 \\
T3.3 & T1571, T1133 \\
T3.4 & T1083, T1005 \\
T3.5 & T1611, T1068 \\
T3.6 & T1199, T1071 \\
T3.7 & T1199, T1556 (partial) \\
T3.8 & T1610, T1525, T1195.002 \\
T3.9 & T1610, T1559, T1041 \\
T3.10 & T1559, T1071, T1078 \\
T5.1 & T1190 (analogue), T1068 \\
T5.2 & T1199, T1570, T1041 \\
T5.3 & T1068, T1554 (analogue) \\
T5.4 & T1068, T1548 \\
\end{longtable}

\begin{center}\rule{0.5\linewidth}{0.5pt}\end{center}

\hypertarget{appendix-d-control-framework-mapping-preliminary}{%
\chapter{Appendix D --- Control-framework mapping
(preliminary)}\label{appendix-d-control-framework-mapping-preliminary}}

For security teams mapping the catalogue to compliance control
frameworks. Definitive mappings require organisation-specific review;
these are first-pass starting points.

\hypertarget{soc-2-trust-services-criteria}{%
\section{SOC 2 (Trust Services
Criteria)}\label{soc-2-trust-services-criteria}}

\begin{longtable}[]{@{}
  >{\raggedright\arraybackslash}p{(\columnwidth - 2\tabcolsep) * \real{0.5000}}
  >{\raggedright\arraybackslash}p{(\columnwidth - 2\tabcolsep) * \real{0.5000}}@{}}
\toprule\noalign{}
\begin{minipage}[b]{\linewidth}\raggedright
T-ID
\end{minipage} & \begin{minipage}[b]{\linewidth}\raggedright
Relevant TSC
\end{minipage} \\
\midrule\noalign{}
\endhead
\bottomrule\noalign{}
\endlastfoot
T1.1, T2.3 & CC6.1, CC6.7 (Logical and Physical Access) \\
T1.2, T1.3, T1.9 & CC8.1 (Change Management); CC7.1 (System Operations:
detection of malicious software) \\
T1.8, T1.6, T1.12 & CC6.6, CC6.7 (boundary controls; transmission of
data) \\
T2.1, T2.5, T2.8 & CC7.2, CC8.1 (system monitoring; change
management) \\
T2.2, T2.4 & CC8.1 (change management); CC7.1 (detection of system
anomalies) \\
T3.2--T3.5 & CC6.6, CC8.1 (boundary controls; change management) \\
T3.6, T3.7 & CC8.1, CC7.1 (change management; detection) \\
T3.8, T3.9, T3.10 & CC6.1, CC6.3 (least-privilege logical access); CC9.2
(vendor and business-partner risk: third-party images and templates) \\
\end{longtable}

\hypertarget{isoiec-270012022-selected-controls}{%
\section{ISO/IEC 27001:2022 (selected
controls)}\label{isoiec-270012022-selected-controls}}

\begin{longtable}[]{@{}
  >{\raggedright\arraybackslash}p{(\columnwidth - 2\tabcolsep) * \real{0.5000}}
  >{\raggedright\arraybackslash}p{(\columnwidth - 2\tabcolsep) * \real{0.5000}}@{}}
\toprule\noalign{}
\begin{minipage}[b]{\linewidth}\raggedright
T-ID
\end{minipage} & \begin{minipage}[b]{\linewidth}\raggedright
Relevant Annex A control
\end{minipage} \\
\midrule\noalign{}
\endhead
\bottomrule\noalign{}
\endlastfoot
T1.1, T2.3 & A.5.10 (Acceptable use); A.8.3 (Information access
restriction); A.8.4 (Access to source code) \\
T1.2, T1.3, T1.9 & A.8.8 (Management of technical vulnerabilities);
A.8.30 (Outsourced development) \\
T1.8, T1.6, T1.12 & A.8.20 (Network security); A.8.21 (Security of
network services) \\
T2.1, T2.5, T2.8 & A.8.31 (Separation of environments); A.8.32 (Change
management) \\
T2.2, T2.4 & A.8.28 (Secure coding); A.8.29 (Security testing in
development and acceptance) \\
T3.2--T3.5 & A.8.22 (Segregation of networks); A.8.23 (Web filtering) \\
T3.6, T3.7 & A.5.7 (Threat intelligence); A.8.28 (Secure coding) \\
T3.8, T3.9, T3.10 & A.8.2 (Privileged access rights); A.8.3 (Information
access restriction); A.5.19, A.5.21 (Supplier relationships; managing
ICT supply chain) \\
\end{longtable}

\hypertarget{nis2-directive-eu-20222555-article-21-measures}{%
\section{NIS2 (Directive (EU) 2022/2555, Article 21
measures)}\label{nis2-directive-eu-20222555-article-21-measures}}

\begin{longtable}[]{@{}
  >{\raggedright\arraybackslash}p{(\columnwidth - 2\tabcolsep) * \real{0.5000}}
  >{\raggedright\arraybackslash}p{(\columnwidth - 2\tabcolsep) * \real{0.5000}}@{}}
\toprule\noalign{}
\begin{minipage}[b]{\linewidth}\raggedright
T-ID
\end{minipage} & \begin{minipage}[b]{\linewidth}\raggedright
Relevant Article 21 measure
\end{minipage} \\
\midrule\noalign{}
\endhead
\bottomrule\noalign{}
\endlastfoot
T1.1--T1.14 & Art. 21(2)(d) supply chain security; 21(2)(e) security in
acquisition, development, maintenance \\
T2.1--T2.8 & Art. 21(2)(a) risk-analysis and information-system security
policies; 21(2)(g) basic cyber-hygiene practices \\
T3.1--T3.7 & Art. 21(2)(d) supply chain security; 21(2)(j) secured
communications \\
T3.8, T3.9, T3.10 & Art. 21(2)(i) access-control policies and least
privilege; 21(2)(d) supply chain security \\
T5.1--T5.4 & Art. 21(2)(e) security in acquisition, development,
maintenance; 21(2)(i) access control and asset management \\
\end{longtable}

\hypertarget{dora-regulation-eu-20222554-for-financial-entities}{%
\section{DORA (Regulation (EU) 2022/2554, for financial
entities)}\label{dora-regulation-eu-20222554-for-financial-entities}}

\begin{longtable}[]{@{}
  >{\raggedright\arraybackslash}p{(\columnwidth - 2\tabcolsep) * \real{0.5000}}
  >{\raggedright\arraybackslash}p{(\columnwidth - 2\tabcolsep) * \real{0.5000}}@{}}
\toprule\noalign{}
\begin{minipage}[b]{\linewidth}\raggedright
T-ID
\end{minipage} & \begin{minipage}[b]{\linewidth}\raggedright
Relevant DORA article
\end{minipage} \\
\midrule\noalign{}
\endhead
\bottomrule\noalign{}
\endlastfoot
T1.1--T1.14 & Art. 9 (ICT risk-management framework); Art. 28
(third-party risk) \\
T2.1--T2.8 & Art. 9; Art. 16, 17 (ICT-related incidents and
reporting) \\
T3.1--T3.7 & Art. 9; Art. 28 (third-party risk) \\
T3.8, T3.9, T3.10 & Art. 9; Art. 28 (third-party risk) \\
T5.1--T5.4 & Art. 9; Art. 16 (ICT-related incidents) \\
\end{longtable}

These mappings are starting points for compliance teams. Specific
control-objective mapping requires organisation-specific analysis.

\begin{center}\rule{0.5\linewidth}{0.5pt}\end{center}

\hypertarget{versioning-and-contribution}{%
\chapter{Versioning and
contribution}\label{versioning-and-contribution}}

Threat identifiers are stable within a published version and across
patch revisions; major-version revisions may renumber. Contributions
follow the structure of existing entries: definition, observed instances
with citations, attack pattern, confinement approach, residuals, ATT\&CK
mapping. Citations should reference public, durable sources.

The catalogue is intended as community infrastructure. Organisations are
encouraged to cite the threat identifiers in their internal policy
documents, and other tools in the user-level-workload-confinement space
are encouraged to adopt them as a shared vocabulary, regardless of
mechanism choices.
